% 同一字体、字号与版心用于中文和法语；复用全学期的字体配置，不修改其源文件。
\input{../../知识整合/排版配置.tex}
\usepackage{tikz}
\providecommand{\blacksquare}{\mathord{\rule{.6em}{.6em}}}
% Preuves p.4 的三幅区间示意，使用 TeX 绘制而非嵌入 PDF 裁图。
\newcommand{\DSintervaldiagram}[1]{%
  \par\begin{center}
  \begin{tikzpicture}[x=1.8cm,y=.5cm,baseline]
    \draw[->] (-1.25,0)--(1.3,0);
    \node at (-1,0) {$]$};\node at (1,0) {$[$};
    \node[below] at (-1,0) {$x-r_x$};\node[below] at (1,0) {$x+r_x$};
    \fill (0,0) circle (1pt);\node[below] at (0,0) {$x$};
    \ifnum#1>1
      \fill (-.25,0) circle (1pt);\node[above] at (-.25,0) {$q_x$};
    \fi
    \ifnum#1=3
      \draw[line width=.7pt] (-.7,.9)--(.2,.9);
      \node at (-.7,.9) {$]$};\node at (.2,.9) {$[$};
      \node[above] at (-.7,.9) {$q_x-\rho_x$};\node[above] at (.2,.9) {$q_x+\rho_x$};
    \fi
  \end{tikzpicture}
  \end{center}%
}
